\documentclass[a4paper,10pt]{article}
\usepackage[utf8]{inputenc}
\input{included_packages} % Required for inserting images

\title{NUR A Hand in 1}
\author{Matthew Kirwan \\
        S4557093}
\date{March 2026}

\begin{document}

\maketitle

\section*{Introduction}
This study examines the use of memory efficient methods with the aim of handling overflow and underflow in the Poisson distribution. Additionally, we implement some linear algebra methods to obtain the 19 degree polynomial from the Vandermonde matrix, and comparing the error against the number of iterations of the LU Decomposition and Nevilles algorithm.

\section*{Problems and Methods}
In this section, the steps taken to solve the given problems are explained. We begin by looking at the initialization of the script.

We import a number of modules for data/path handling, graphing, and timing purposes. Here we automatically grab the directory the script is in, and navigate around this path to access the directory structure surrounding it.

\begin{lstlisting}[language=Python]
import os
import time
import math
import numpy as np
import pandas as pd

from functions import interpolation
from functions import math_functions
from functions import linear_algebra as la

import matplotlib.pyplot as plt

# get the folder where the script is located
script_dir = os.path.dirname(__file__)
data_dir = os.path.join(script_dir, "../data")
plots_dir = os.path.join(script_dir, "../plots")

os.makedirs(plots_dir, exist_ok=True)  # just in case
\end{lstlisting}

In the section above we import 3 files which will be used in solving the problems proposed, 'interpolation', 'math\_functions', and 'linear\_algebra'. We will explore these further upon their usage.

\subsection*{Question 1: Poisson Distribution}
The problem proposed here is focused around handling over/underflow through the calculation of the Poisson distribution, defined in eq. \ref{eq:poisson}

\begin{equation}
    P_\lambda(K) = \frac{\lambda^ke^{-\lambda}}{k!}
    \label{eq:poisson}
\end{equation}

Here we are restricted to the use of 32 bit variables, so we utilise the numpy.float32 and numpy.int32 functions to ensure this restriction is met. In order to test the effectiveness of our algorithm, we are given a number of values which are defined in table \ref{tab:lam_and_k}. We are also expected to output to at least six significant figures.

\begin{table}[]
    \centering
    \begin{tabular}{|c|c|}
        \hline
        $\lambda$ & k \\
        \hline
        1 & 0 \\
        5 & 10 \\
        3 & 21 \\
        2.6 & 40 \\
        100 & 5 \\
        101 & 200 \\
        \hline
    \end{tabular}
    \caption{$\lambda$ and k values given}
    \label{tab:lam_and_k}
\end{table}

We begin by defining the test data in the question\_1\_values numpy array, ensuring the data is stored in numpy.float32 bits. We then then define an empty list for appending our results to later and begin looping through the test variables. In order to handle both the integer and float $\lambda$ values, we load them in as floats initially, and check if the float value is an integer, if it is we convert this $\lambda$ to an numpy.int32 variable, otherwise we do nothing. We then ensure that k is an numpy.int32 variable as the factorial can only ever handle integers.

We can now safely pass the values through to the Poisson distribution defined in the math\_functions, this returns a single scalar value which we then append to the results list as an array and then convert this to a pandas DataFrame for pretty printing.

\begin{lstlisting}[language=Python]
def question_1() -> pd.DataFrame:
    question_1_values = np.array([[1,0],[5,10],[3,21],[2.6,40],[100,5],[101,200]],dtype=np.float32)

    results = []
    for values in question_1_values:

        #defining how the data is stored in 32 bits
        lam_raw = np.float32(values[0])
        if lam_raw.is_integer():
            lam = np.int32(lam_raw)
        else:
            lam = lam_raw
        k = np.int32(values[1])

        result = math_functions.poisson_distribution(lam,k)
        results.append([lam,k,result])
    results = pd.DataFrame(results,columns=['lambda','k','Poisson result'])
    return results
\end{lstlisting}

When we pass the variables to the Poisson distribution function, we check if k is 0, in the case it is, we calculate the distribution as it is written in eq. \ref{eq:poisson}. This is truly just an artifact of early prototyping, as the other condition can handle the log(0) due to the fact that there is no division. In the other case, we apply a logarithm to the entire distribution, which results in the equation to be re-written as
eq. \ref{eq: log_poisson}. Once the values are obtained, they are returned as numpy.float32 variables.

\begin{equation}
    Log(P_{(\lambda,k)}) = kLog(\lambda) - \lambda - Log(k!)
    \label{eq: log_poisson}
\end{equation}

\newpage
\begin{lstlisting}[language=Python]
def poisson_distribution(lam:np.float32|np.int32,
                         k:np.int32
                         ) -> np.int32:

    if k == 0:
        p = ((lam ** k) * np.exp(-lam)) / (factorial(k))
        return np.float32(p)
    else:
        p = k * math.log(lam) - lam - log_factorial(k)
        return np.exp(np.float32(p))
\end{lstlisting}

We can encounter overflow easily if k is double digits, to handle this, we implement a logarithmic factorial function. Regardless, the factorial function is a recursive function that returns the product of the previous factorials once 1 is reached. These values are returned in numpy.float32 variables so that they can be returned as an exponent.

The log factorial function is not recursive but a for loop which obtains the summation of the Log of each value until the k integer is reached. As we initialize the sum with 0, and don't handle the log of 0, this function can easily handle the k = 0 case. As we require the log of the factorial, this value is returned as a numpy.float32 variable.

\begin{lstlisting}[language=Python]
#tried to use this initially but failed at k = 40 due to overflow
#tried return log(product) but realised that was a recursive log return so...
def factorial(x:np.int32
              ) -> np.float32:

    if x > 1:
        return np.float32(x * factorial(x-1))
    else:
        return np.float32(1)

#solved overflow issue
def log_factorial(x:np.int32
                  ) -> np.float32:

    fac_sum = 0.0
    for i in range(1, x + 1):
        fac_sum += math.log(i)

    return np.float32(fac_sum)
\end{lstlisting}

\subsection*{Question 2: Vandermonde Matrix}
In this section we are first challenged with obtaining the Vandermonde Matrix \textbf{V}, which is defined by eq. \ref{eq: vandermonde}. Where $x_i$ is 20 points defined in the vandermonde.txt file. Secondly we are asked to obtain the 19 degree polynomial for these 20 points through the use of LU decomposition and 19 degree polynomial interpolation, and timing the time required for each process.

\begin{equation}
    V_{i,j} = x_i^j
    \label{eq: vandermonde}
\end{equation}

I've removed some sections from this function for ease of reading (ie. plotting, saving) and the return declaration is incorrect here, as we return a tuple of multiple arrays. We start this section by importing the data from the vandermonde.txt file. Splitting the data and reorganising the data in to more usable formats, and initialise our storage.

We then calculate the vandermonde matrix with the math\_functions function get\_vandermonde\_matrix(x). Once we have this we can perform our LU decompositions over 1 iteration and 10 iterations with the functions LU\_n\_iter and can time this process with the perf\_counter, we opted out of the usage of timeit as we require the outputs of these functions.

We then perform a timing analysis on the nevilles algorithm implementation by parsing through the 1001 x points and interpolating their values. This function utilises the M variable which is the amount of points around the requested data point (found through bisection) that will be considered during the interpolation, setting M = 20, means all points are considered at all interpolations resulting in a 19 degree polynomial fit. Setting this value lower to 3/5/8 allows for a much smoother approximation as not all points are considered, only neighbouring points.

After this we interpolate the values for all 1001 points through linear/lagrange/Nevilles algorithm interpolation. These values are saved and later plotted. Subsequently, we re-interpolate values for the known x values and obtain the difference between the known value $y_i$ and the interpolated value y(x)

\begin{lstlisting}[language=Python]
def question_2() -> tuple[pd.DataFrame]:
    data = import_data()
    x,y = np.array(data['x']),np.array(data['y'])
    x_points_to_interpolate = np.linspace(x[0],x[-1],1001)
    M=20 #change to 3/5/8 for smooth function
    nevilles    = []
    lagrange    = []
    linear      = []
    x_points = []
    x_y = np.array([x, y])

    vandermonde = math_functions.get_vandermonde_matrix(x)

    start = time.perf_counter()
    c, ys = LU_1_iter(vandermonde, x, y)
    end = time.perf_counter()
    elapsed_1_iter = end - start

    start = time.perf_counter()
    c_1, yyy = LU_10_iter(vandermonde, x, y)
    end = time.perf_counter()
    elapsed_10_iter = end - start

    start = time.perf_counter()
    for point_index in range(0,len(x_points_to_interpolate)):

        point_request = x_points_to_interpolate[point_index]
        interpolated_values_nev     = interpolation.nevilles_algorithm(point_request,M,x_y)
    end = time.perf_counter()
    elapsed_nev = end - start

    #need to split this up for timing purposes
    for point_index in range(0,len(x_points_to_interpolate)):
        point_request = x_points_to_interpolate[point_index]
        interpolated_values_linear  = interpolation.linear_interpolation(point_request,x_y)
        interpolated_values_nev     = interpolation.nevilles_algorithm(point_request,M,x_y)
        interpolated_values_lag     = interpolation.lagrange_polynomial(point_request,19,x_y)
        linear.append(interpolated_values_linear[1])
        nevilles.append(interpolated_values_nev[1])
        lagrange.append(interpolated_values_lag[1])
        x_points.append(x_points_to_interpolate[point_index])

    y_check=[]
    for point_request in x:
        interpolated_values_nev = interpolation.nevilles_algorithm(point_request,M,x_y)
        y_check.append(interpolated_values_nev[1])
    diff_nev = abs(y_check - y)

    y_check = []
    c_1, y_check = LU_1_iter(vandermonde, x, y,20)
    diff_1st = abs(y_check - y)

    y_check = []
    c_1, y_check = LU_10_iter(vandermonde, x, y,20)
    diff_10th = abs(y_check - y)

\end{lstlisting}

\begin{lstlisting}[language=Python]

\end{lstlisting}

\section*{Results}
\subsection*{Question 1: Poisson Distribution}

\subsection*{Question 2: Vandermonde Matrix}

\end{document}
